\chapter{Van gen tot eiwit}\label{ch:g11:gene-expression}

In 1961 voedde een biochemicus een celvrij extract van bacteriën met een
kunstmatig \omterm{def:g11:gene-expression:transcription}{RNA} dat uit niets dan de letter U bestond, telkens herhaald.
Het extract maakte een \omterm{def:g11:gene-expression:protein}{eiwit} dat uit niets dan één \omterm{def:g10:chemistry-of-life:families}{aminozuur} bestond,
fenylalanine, telkens herhaald. Met die ene proef was het eerste woord
van de \omterm{def:g11:gene-expression:code}{genetische code} gelezen: UUU betekent fenylalanine. Binnen vijf
jaar waren alle vierenzestig woorden bekend, en het bleken dezelfde
woorden te zijn bij een bacterie, een tarweplant en een mens. Dit
hoofdstuk volgt de weg van de nucleotidenvolgorde van een \omterm{def:g10:universal-dna:gene}{gen} naar de
aminozuurvolgorde van een \omterm{def:g11:gene-expression:protein}{eiwit} --- de weg waarlangs elk kenmerk uit
\cref{ch:g11:mutations} loopt.

\section{Genen maken eiwitten}

\begin{proposition}[Eén gen, één eiwit]\label{prop:g11:gene-expression:onegene}
Een \omterm{def:g10:universal-dna:gene}{gen} komt tot uiting wanneer de \omterm{def:g10:cells-common-unit:cell}{cel} zijn volgorde gebruikt om het
\omterm{def:g11:gene-expression:protein}{eiwit} te maken waarvoor het codeert. De meeste kenmerken hangen van
\omterm{def:g10:chemistry-of-life:families}{eiwitten} af --- \omterm{def:g10:cell-metabolism:metabolism}{enzymen}, structuurvezels, transporteurs, receptoren ---
en een \omterm{def:g11:mutations:mutation}{mutatie} verandert een kenmerk door het \omterm{def:g11:gene-expression:protein}{eiwit} te veranderen dat
haar \omterm{def:g10:universal-dna:gene}{gen} voortbrengt.
\end{proposition}
\begin{proof}[Bewijsmateriaal]
Beadle en Tatum (1941) bestraalden sporen van een broodschimmel en
verzamelden mutanten die niet meer op een minimaal medium konden groeien
tenzij er één bepaalde stof werd toegevoegd, zoals het \omterm{def:g10:chemistry-of-life:families}{aminozuur}
arginine. Elke mutant miste één \omterm{def:g10:cell-metabolism:metabolism}{enzym} uit de reactieketen die die stof
maakt; elk gebrek werd als één enkel \omterm{def:g10:universal-dna:gene}{gen} geërfd; en mutanten die op
verschillende stappen van dezelfde keten vastliepen, droegen \omterm{def:g11:mutations:mutation}{mutaties} in
verschillende \omterm{def:g10:universal-dna:gene}{genen}. Eén \omterm{def:g10:universal-dna:gene}{gen}, één \omterm{def:g10:cell-metabolism:metabolism}{enzym} --- en, zoals later werk
veralgemeende, één \omterm{def:g10:universal-dna:gene}{gen}, één polypeptideketen.
\end{proof}

\begin{definition}[Eiwit, aminozuurvolgorde]\label{def:g11:gene-expression:protein}
Een \emph{eiwit}\index{eiwit} is een keten van \omterm{def:g10:chemistry-of-life:families}{aminozuren}, van enkele
tientallen tot enkele duizenden lang, gekozen uit een verzameling van
twintig typen. De \emph{volgorde} --- de opeenvolging van de
\omterm{def:g10:chemistry-of-life:families}{aminozuren} --- is wat het \omterm{def:g10:universal-dna:gene}{gen} vastlegt; eenmaal gemaakt vouwt de keten
zich tot een bepaalde driedimensionale vorm die door die volgorde is
bepaald, en de vorm bepaalt de functie. Verander één \omterm{def:g10:chemistry-of-life:families}{aminozuur} en de
vouwing, en dus de functie, kan veranderen.
\end{definition}

\section{Transcriptie: het gen gekopieerd naar RNA}

\begin{definition}[Boodschapper-RNA en transcriptie]\label{def:g11:gene-expression:transcription}
\emph{RNA}\index{RNA} (ribonucleïnezuur) is een enkelstrengs keten van
\omterm{def:g10:universal-dna:nucleotide}{nucleotiden} zoals die van het \omterm{def:g10:universal-dna:information}{DNA}, alleen is de suiker ribose en vervangt
de base uracil (U) de thymine (T). De \emph{transcriptie}\index{transcriptie}
is het overschrijven van de volgorde van een \omterm{def:g10:universal-dna:gene}{gen} naar een molecuul
\emph{boodschapper-RNA}\index{boodschapper-RNA} (mRNA): het \omterm{def:g10:cell-metabolism:metabolism}{enzym}
\emph{RNA-polymerase} opent de \omterm{prop:g10:universal-dna:helix}{dubbele helix} langs het \omterm{def:g10:universal-dna:gene}{gen}, leest één
streng (de \emph{matrijsstreng}) en zet daar het complementaire RNA op
elkaar, U tegenover A, A tegenover T, G tegenover C en C tegenover G. Het
mRNA heeft de volgorde van de andere streng van het \omterm{def:g10:universal-dna:gene}{gen}, met U in plaats
van T; het verlaat de kern naar het \omterm{def:g10:cells-common-unit:cell}{cytoplasma}.
\end{definition}

\begin{omfigure}
\begin{tikzpicture}[font=\small, >=Stealth, line cap=round]
  % DNA strands with bubble
  \draw[very thick, omDef] (0,0.5) -- (2.5,0.5) .. controls (3.2,0.5) and (3.4,1.3) .. (4.3,1.3) -- (6.7,1.3) .. controls (7.6,1.3) and (7.8,0.5) .. (8.5,0.5) -- (11,0.5);
  \draw[very thick, omDef] (0,-0.5) -- (2.5,-0.5) .. controls (3.2,-0.5) and (3.4,-1.3) .. (4.3,-1.3) -- (6.7,-1.3) .. controls (7.6,-1.3) and (7.8,-0.5) .. (8.5,-0.5) -- (11,-0.5);
  \foreach \x in {0.3,0.7,...,2.3} \draw[gray] (\x,0.5) -- (\x,-0.5);
  \foreach \x in {8.7,9.1,...,10.9} \draw[gray] (\x,0.5) -- (\x,-0.5);
  % template strand letters (bottom strand in bubble)
  \foreach \x/\l in {4.5/T, 4.9/A, 5.3/C, 5.7/G, 6.1/G, 6.5/A} \node[font=\scriptsize, omDef, anchor=north] at (\x,-1.32) {\l};
  % mRNA
  \draw[very thick, omThm] (4.4,-0.9) -- (6.6,-0.9) .. controls (7.4,-0.9) and (7.2,0.1) .. (8.0,0.6) .. controls (8.6,1.0) and (9.4,1.4) .. (10.6,1.7);
  \foreach \x/\l in {4.5/A, 4.9/U, 5.3/G, 5.7/C, 6.1/C, 6.5/U} \node[font=\scriptsize, omThm, anchor=south] at (\x,-0.85) {\l};
  \foreach \x in {4.5,4.9,...,6.5} \draw[gray, densely dotted] (\x,-0.95) -- (\x,-1.28);
  % polymerase
  \draw[thick, fill=omProp!30, rounded corners=8pt] (6.95,-1.6) rectangle (8.65,0.2);
  \node[align=center] at (7.8,-0.7) {RNA-\\ polymerase};
  \draw[->, thick] (8.7,-1.9) -- (9.8,-1.9) node[right] {schuift langs het gen};
  \node[anchor=east, omDef, align=right] at (-0.15,0) {DNA van\\ het gen};
  \node[anchor=west, omThm, align=left] at (10.7,1.7) {mRNA, dat\\ vrijkomt};
  \node[anchor=north, omDef] at (5.5,-1.75) {matrijsstreng};
\end{tikzpicture}

{\small De \omterm{def:g11:gene-expression:transcription}{transcriptie}. Het RNA-polymerase opent de helix, leest de
matrijsstreng en bouwt daarop een \omterm{def:g11:gene-expression:transcription}{boodschapper-RNA} dat er complementair
aan is (U in plaats van T). Achter het \omterm{def:g10:cell-metabolism:metabolism}{enzym} sluit de helix zich weer;
het \omterm{def:g11:gene-expression:transcription}{RNA} laat los.}
\end{omfigure}

\begin{example}[Een transcript]\label{ex:g11:gene-expression:transcript}
De matrijsstreng \texttt{3'-TAC GGA CTT ATC-5'} geeft het mRNA
\texttt{5'-AUG CCU GAA UAG-3'}. De andere streng van het \omterm{def:g10:universal-dna:gene}{gen} luidt
\texttt{5'-ATG CCT GAA TAG-3'}: het mRNA is de volgorde van die streng
met U in plaats van T, en daarom worden genvolgordes bij afspraak als
die streng geschreven, de \emph{coderende streng}.
\end{example}

\section{De genetische code}

\begin{definition}[Codon en de genetische code]\label{def:g11:gene-expression:code}
Het mRNA wordt vanaf een vast beginpunt in opeenvolgende groepjes van
drie \omterm{def:g10:universal-dna:nucleotide}{nucleotiden} gelezen, de \emph{codons}\index{codon}. De
\emph{genetische code}\index{genetische code} is het verband tussen de 64
mogelijke codons en de 20 \omterm{def:g10:chemistry-of-life:families}{aminozuren}: 61 codons wijzen een \omterm{def:g10:chemistry-of-life:families}{aminozuur} aan,
drie (UAA, UAG, UGA) zijn \emph{stopsignalen} die de keten beëindigen, en
AUG, dat methionine aanwijst, dient tevens als \emph{startcodon}. De code
is \emph{ontaard} --- de meeste \omterm{def:g10:chemistry-of-life:families}{aminozuren} hebben verscheidene codons ---
en \emph{universeel}: dezelfde tabel geldt voor elk bekend organisme, op
zeldzame kleine varianten na.
\end{definition}

\begin{omfigure}
\begin{tabular}{c|cccc|c}
\multicolumn{1}{c}{} & \multicolumn{4}{c}{\small tweede letter} & \multicolumn{1}{c}{} \\
{\small eerste} & U & C & A & G & {\small derde} \\ \hline
U & UUU Phe & UCU Ser & UAU Tyr & UGU Cys & U \\
 & UUC Phe & UCC Ser & UAC Tyr & UGC Cys & C \\
 & UUA Leu & UCA Ser & UAA \textbf{stop} & UGA \textbf{stop} & A \\
 & UUG Leu & UCG Ser & UAG \textbf{stop} & UGG Trp & G \\ \hline
C & CUU Leu & CCU Pro & CAU His & CGU Arg & U \\
 & CUC Leu & CCC Pro & CAC His & CGC Arg & C \\
 & CUA Leu & CCA Pro & CAA Gln & CGA Arg & A \\
 & CUG Leu & CCG Pro & CAG Gln & CGG Arg & G \\ \hline
A & AUU Ile & ACU Thr & AAU Asn & AGU Ser & U \\
 & AUC Ile & ACC Thr & AAC Asn & AGC Ser & C \\
 & AUA Ile & ACA Thr & AAA Lys & AGA Arg & A \\
 & AUG \textbf{Met} & ACG Thr & AAG Lys & AGG Arg & G \\ \hline
G & GUU Val & GCU Ala & GAU Asp & GGU Gly & U \\
 & GUC Val & GCC Ala & GAC Asp & GGC Gly & C \\
 & GUA Val & GCA Ala & GAA Glu & GGA Gly & A \\
 & GUG Val & GCG Ala & GAG Glu & GGG Gly & G \\
\end{tabular}

{\small De \omterm{def:g11:gene-expression:code}{genetische code}, gelezen op het mRNA. Elk \omterm{def:g11:gene-expression:code}{codon} wordt gevonden
via zijn eerste letter (rij), zijn tweede letter (kolom) en zijn derde
letter (regel binnen het blok). AUG is zowel methionine als het
startsignaal; drie \omterm{def:g11:gene-expression:code}{codons} zijn stops.}
\end{omfigure}

\begin{proposition}[Hoe de code werd gelezen]\label{prop:g11:gene-expression:readcode}
De betekenis van elk \omterm{def:g11:gene-expression:code}{codon} is proefondervindelijk vastgesteld.
\end{proposition}
\begin{proof}[Bewijsmateriaal]
Nirenberg en Matthaei (1961) voegden kunstmatig \omterm{def:g11:gene-expression:transcription}{RNA} van één herhaalde
\omterm{def:g10:universal-dna:nucleotide}{nucleotide} toe aan een celvrij extract met \omterm{def:g10:cells-common-unit:organelle}{ribosomen}, \omterm{def:g10:chemistry-of-life:families}{aminozuren} en de
\omterm{def:g10:cell-metabolism:metabolism}{enzymen} van de eiwitsynthese: poly-U leverde een keten fenylalanines op,
poly-C een keten prolines, poly-A een keten lysines. Kunstmatige \omterm{def:g11:gene-expression:transcription}{RNA}'s
van herhaalde paren en drietallen, en later het binden van losse
drietallen aan \omterm{def:g10:cells-common-unit:organelle}{ribosomen}, wezen de resterende \omterm{def:g11:gene-expression:code}{codons} in 1966 toe. De
lengte van drie letters was uit de rekenkunde afgeleid (twee letters
geven slechts 16 combinaties, te weinig voor 20 \omterm{def:g10:chemistry-of-life:families}{aminozuren}; drie geven er
64) en door \omterm{def:g11:mutations:mutation}{mutaties} bevestigd: één of twee \omterm{def:g10:universal-dna:nucleotide}{nucleotiden} in een \omterm{def:g10:universal-dna:gene}{gen}
inbrengen vernietigt zijn \omterm{def:g11:gene-expression:protein}{eiwit}, drie inbrengen herstelt een vrijwel
normaal \omterm{def:g11:gene-expression:protein}{eiwit}.
\end{proof}

\begin{method}[Een volgorde vertalen]\label{met:g11:gene-expression:translate}
\begin{enumerate}
  \item Schrijf de coderende streng op (of transcribeer de
        matrijsstreng): vervang T door U om het mRNA te krijgen.
  \item Zoek de eerste AUG: die is het begin, en daarmee ligt het
        leesraam vast.
  \item Knip het mRNA vanaf die AUG in opeenvolgende drietallen en zoek
        elk drietal in de tabel op.
  \item Stop bij het eerste stopcodon; de \omterm{def:g10:chemistry-of-life:families}{aminozuren} die je op volgorde
        hebt opgeschreven zijn de volgorde van het \omterm{def:g11:gene-expression:protein}{eiwit}.
  \item Doe hetzelfde voor een mutant en vergelijk: hetzelfde \omterm{def:g11:gene-expression:protein}{eiwit}
        (stil), één \omterm{def:g10:chemistry-of-life:families}{aminozuur} veranderd (missense), een voortijdige stop
        (nonsense), alles na de \omterm{def:g11:mutations:mutation}{mutatie} veranderd (leesraamverschuiving).
\end{enumerate}
\end{method}

\begin{example}[Vier mutaties, vier uitkomsten]\label{ex:g11:gene-expression:outcomes}
Coderende streng \texttt{ATG CCT GAA TTC TAG}: mRNA
\texttt{AUG CCU GAA UUC UAG}, \omterm{def:g11:gene-expression:protein}{eiwit} Met--Pro--Glu--Phe.
\begin{itemize}
  \item \texttt{ATG CC\textbf{C} GAA TTC TAG}: CCC is nog steeds Pro ---
        \emph{stil}.
  \item \texttt{ATG CCT G\textbf{C}A TTC TAG}: GCA is Ala ---
        Met--Pro--Ala--Phe, \emph{missense}.
  \item \texttt{ATG CCT \textbf{T}AA TTC TAG}: UAA is een stop ---
        Met--Pro, \emph{nonsense}, een afgekapt \omterm{def:g11:gene-expression:protein}{eiwit}.
  \item \texttt{ATG CC\textbf{A}T GAA TTC TAG}: een ingevoegde A
        verschuift het raam: \texttt{AUG CCA UGA}\dots{} Met--Pro--stop
        --- een \emph{leesraamverschuiving}.
\end{itemize}
\end{example}

\section{Translatie: de boodschap gelezen}

\begin{proposition}[Translatie]\label{prop:g11:gene-expression:translation}
De \emph{translatie}\index{translatie} vindt plaats op de
\emph{\omterm{def:g10:cells-common-unit:organelle}{ribosomen}}, in het \omterm{def:g10:cells-common-unit:cell}{cytoplasma}. Een \omterm{def:g10:cells-common-unit:organelle}{ribosoom} bindt het mRNA bij zijn
startcodon en schuift er \omterm{def:g11:gene-expression:code}{codon} voor \omterm{def:g11:gene-expression:code}{codon} langs. Bij elk \omterm{def:g11:gene-expression:code}{codon} hoort een
\emph{\omterm{def:g11:gene-expression:transcription}{transport-RNA}} (tRNA), een klein \omterm{def:g11:gene-expression:transcription}{RNA} dat aan het ene uiteinde het
\emph{anticodon} van drie \omterm{def:g10:universal-dna:nucleotide}{nucleotiden} draagt dat complementair is met het
\omterm{def:g11:gene-expression:code}{codon}, en aan het andere uiteinde het bijbehorende \omterm{def:g10:chemistry-of-life:families}{aminozuur}. Het
\omterm{def:g10:cells-common-unit:organelle}{ribosoom} koppelt elk aangevoerd \omterm{def:g10:chemistry-of-life:families}{aminozuur} aan de groeiende keten en laat
het lege tRNA los; bij een stopcodon laat het de voltooide keten los, die
zich vouwt. Verscheidene \omterm{def:g10:cells-common-unit:organelle}{ribosomen} lezen één mRNA na elkaar, en één mRNA
levert honderden kopieën van het \omterm{def:g11:gene-expression:protein}{eiwit} op voordat het wordt afgebroken.
\end{proposition}
\admitted

\begin{omfigure}
\begin{tikzpicture}[font=\small, >=Stealth, line cap=round]
  % mRNA
  \draw[very thick, omThm] (0,0) -- (11,0);
  \foreach \x/\l in {0.4/A,0.8/U,1.2/G, 1.7/C,2.1/C,2.5/U, 3.0/G,3.4/A,3.8/A, 4.3/U,4.7/U,5.1/C, 5.6/G,6.0/G,6.4/U, 6.9/U,7.3/A,7.7/G}
    \node[font=\scriptsize, omThm, anchor=north] at (\x,-0.05) {\l};
  \node[anchor=east, omThm] at (-0.15,0) {mRNA};
  % ribosome (two lobes) over codons 3 and 4
  \draw[thick, fill=omProp!25, rounded corners=10pt] (2.7,0.15) rectangle (5.5,1.4);
  \draw[thick, fill=omProp!35, rounded corners=12pt] (2.5,1.4) .. controls (2.5,3.0) and (5.7,3.0) .. (5.7,1.4) -- cycle;
  \node at (4.1,2.1) {ribosoom};
  % tRNAs
  \draw[very thick, omMeth!80!black] (3.4,0.55) -- (3.4,1.35);
  \node[font=\scriptsize, omMeth!80!black, anchor=south] at (3.4,0.15) {CUU};
  \draw[very thick, omMeth!80!black] (4.7,0.55) -- (4.7,1.35);
  \node[font=\scriptsize, omMeth!80!black, anchor=south] at (4.7,0.15) {AAG};
  % amino acids chain
  \foreach \i/\x in {0/3.4, 1/4.7} \fill[omDef!70] (\x,1.5) circle (4pt);
  \foreach \x in {2.6, 1.9, 1.2} \fill[omDef!70] (\x,1.9) circle (4pt);
  \draw[thick, omDef!70] (1.2,1.9) -- (2.6,1.9) -- (3.4,1.5) -- (4.7,1.5);
  \node[anchor=east, omDef!80!black, align=right] at (0.9,1.9) {groeiend\\ eiwit};
  % incoming tRNA
  \draw[very thick, omMeth!80!black] (7.0,1.3) -- (7.0,2.1);
  \fill[omDef!70] (7.0,2.25) circle (4pt);
  \node[font=\scriptsize, omMeth!80!black, anchor=west] at (7.15,1.35) {CCA};
  \node[anchor=west, align=left, font=\scriptsize] at (7.3,1.7) {volgend tRNA, met zijn\\ aminozuur, komt aan};
  \draw[->, thick] (7.0,1.1) -- (6.1,0.75);
  \draw[->, thick] (5.7,-0.9) -- (7.2,-0.9) node[right] {het ribosoom schuift op};
  \node[anchor=north, font=\scriptsize] at (4.05,-0.45) {codon};
  \draw (2.85,-0.35) -- (2.85,-0.45) -- (3.95,-0.45) -- (3.95,-0.35);
\end{tikzpicture}

{\small De \omterm{prop:g11:gene-expression:translation}{translatie}. Het \omterm{def:g10:cells-common-unit:organelle}{ribosoom} houdt twee \omterm{def:g11:gene-expression:code}{codons} vast; een tRNA
waarvan het anticodon bij een \omterm{def:g11:gene-expression:code}{codon} past, brengt zijn \omterm{def:g10:chemistry-of-life:families}{aminozuur}, de keten
wordt op de nieuwkomer overgezet, en het \omterm{def:g10:cells-common-unit:organelle}{ribosoom} schuift één \omterm{def:g11:gene-expression:code}{codon} op.
Bij een stopcodon wordt de keten losgelaten.}
\end{omfigure}

\begin{example}[Snelheid en opbrengst]\label{ex:g11:gene-expression:speed}
Een \omterm{def:g10:cells-common-unit:organelle}{ribosoom} voegt 5 tot 20 \omterm{def:g10:chemistry-of-life:families}{aminozuren} per seconde toe: een \omterm{def:g11:gene-expression:protein}{eiwit} van 300
\omterm{def:g10:chemistry-of-life:families}{aminozuren} kost minder dan een minuut. \omterm{def:g10:cells-common-unit:organelle}{Ribosomen} volgen elkaar op een
mRNA met ongeveer 80 \omterm{def:g10:universal-dna:nucleotide}{nucleotiden} ertussen, dus draagt een boodschap van
900 \omterm{def:g10:universal-dna:nucleotide}{nucleotiden} er een tiental tegelijk; over haar leven van enkele uren
levert één mRNA duizend eiwitmoleculen op. Een \omterm{def:g10:cells-common-unit:cell}{cel} van \emph{E.~coli}
maakt zo'n $\num{20000}$ eiwitmoleculen per seconde; de voorloper van een
rode bloedcel zo'n $\num{5e8}$ hemoglobinemoleculen in haar leven.
\end{example}

\section{Eén gen, verscheidene boodschappen}

\begin{proposition}[Rijping van de boodschap bij eukaryoten]\label{prop:g11:gene-expression:splicing}
Bij eukaryote \omterm{def:g10:cells-common-unit:cell}{cellen} wordt de volgorde van een \omterm{def:g10:universal-dna:gene}{gen} onderbroken door
\emph{intronen}\index{intron}, stukken die wel worden overgeschreven maar
daarna uit het \omterm{def:g11:gene-expression:transcription}{RNA} worden geknipt voordat het de kern verlaat; de stukken
die behouden blijven en aaneen worden gekoppeld, de \emph{exonen}, vormen
het rijpe mRNA. Veel \omterm{def:g10:universal-dna:gene}{genen} kunnen op meer dan één manier worden geknipt,
waarbij verschillende stellen exonen worden behouden (\emph{alternatieve
splitsing}), zodat één \omterm{def:g10:universal-dna:gene}{gen} verscheidene verschillende mRNA's en
verscheidene verwante \omterm{def:g10:chemistry-of-life:families}{eiwitten} oplevert: de $\num{20000}$ menselijke
\omterm{def:g10:universal-dna:gene}{genen} coderen voor ruim $\num{100000}$ \omterm{def:g10:chemistry-of-life:families}{eiwitten}.
\end{proposition}
\admitted

\begin{omfigure}
\begin{tikzpicture}[font=\small, >=Stealth]
  \tikzset{ex/.style={draw, fill=omDef!40, minimum height=0.45cm, minimum width=1.0cm, inner sep=1pt},
           intr/.style={draw, fill=gray!20, minimum height=0.45cm, minimum width=0.7cm, inner sep=1pt}}
  \node[anchor=east] at (-0.2,0) {gen (DNA)};
  \node[ex] at (0.5,0) {1}; \node[intr] at (1.35,0) {}; \node[ex] at (2.2,0) {2}; \node[intr] at (3.05,0) {}; \node[ex] at (3.9,0) {3}; \node[intr] at (4.75,0) {}; \node[ex] at (5.6,0) {4};
  \node[anchor=west, font=\scriptsize] at (6.2,0) {exonen (blauw), intronen (grijs)};
  \draw[->, thick] (3.0,-0.35) -- (3.0,-0.95) node[midway, right] {transcriptie};
  \node[anchor=east] at (-0.2,-1.3) {pre-mRNA (onrijp)};
  \node[ex, fill=omThm!30] at (0.5,-1.3) {1}; \node[intr] at (1.35,-1.3) {}; \node[ex, fill=omThm!30] at (2.2,-1.3) {2}; \node[intr] at (3.05,-1.3) {}; \node[ex, fill=omThm!30] at (3.9,-1.3) {3}; \node[intr] at (4.75,-1.3) {}; \node[ex, fill=omThm!30] at (5.6,-1.3) {4};
  \draw[->, thick] (2.0,-1.65) -- (1.0,-2.35) node[midway, left] {splitsing};
  \draw[->, thick] (4.0,-1.65) -- (6.0,-2.35) node[midway, right] {alternatieve splitsing};
  \node[anchor=east] at (-0.6,-2.7) {mRNA A};
  \node[ex, fill=omThm!30] at (0.0,-2.7) {1}; \node[ex, fill=omThm!30] at (1.0,-2.7) {2}; \node[ex, fill=omThm!30] at (2.0,-2.7) {3}; \node[ex, fill=omThm!30] at (3.0,-2.7) {4};
  \node[anchor=west] at (3.6,-2.7) {eiwit A};
  \node[anchor=east] at (5.4,-2.7) {};
  \node[ex, fill=omThm!30] at (6.0,-2.7) {1}; \node[ex, fill=omThm!30] at (7.0,-2.7) {2}; \node[ex, fill=omThm!30] at (8.0,-2.7) {4};
  \node[anchor=west] at (8.6,-2.7) {eiwit B};
\end{tikzpicture}

{\small Een eukaryoot \omterm{def:g10:universal-dna:gene}{gen} wordt in zijn geheel overgeschreven, waarna de
\omterm{prop:g11:gene-expression:splicing}{intronen} worden verwijderd. Alle exonen behouden, of er één weglaten,
geeft twee mRNA's en twee \omterm{def:g10:chemistry-of-life:families}{eiwitten} uit hetzelfde \omterm{def:g10:universal-dna:gene}{gen}.}
\end{omfigure}

\begin{remark}[De stroom van de informatie]\label{rem:g11:gene-expression:flow}
\omterm{def:g10:universal-dna:information}{DNA} wordt overgeschreven naar \omterm{def:g11:gene-expression:transcription}{RNA}, \omterm{def:g11:gene-expression:transcription}{RNA} wordt vertaald naar \omterm{def:g11:gene-expression:protein}{eiwit}, en het
\omterm{def:g11:gene-expression:protein}{eiwit} maakt het kenmerk: de informatie stroomt maar één kant op. Een
veranderd \omterm{def:g11:gene-expression:protein}{eiwit} herschrijft nooit het \omterm{def:g10:universal-dna:gene}{gen}, en daarom wordt een leven lang
trainen of verbranden in de zon niet geërfd, en veranderen alleen
\omterm{def:g11:mutations:mutation}{mutaties} van het \omterm{def:g10:universal-dna:information}{DNA} (\cref{ch:g11:mutations}) wat de volgende generatie
ontvangt. Dezelfde drie stappen lopen in elke \omterm{def:g10:cells-common-unit:cell}{cel} van elk organisme, in
dezelfde code: de universaliteit van \cref{ch:g10:universal-dna}, nu tot
op het laatste woord.
\end{remark}

\section{Oefeningen}

\begin{exercise}[$\star$]\label{exo:g11:gene-expression:1}
Geef drie verschillen tussen \omterm{def:g10:universal-dna:information}{DNA} en \omterm{def:g11:gene-expression:transcription}{RNA}.
\end{exercise}

\begin{exercise}[$\star$]\label{exo:g11:gene-expression:2}
Transcribeer de matrijsstreng \texttt{3'-TAC CGA AAA ATT-5'} naar mRNA.
\end{exercise}

\begin{exercise}[$\star$]\label{exo:g11:gene-expression:3}
Vertaal het mRNA \texttt{AUG GGC AAA UGU UAA} met behulp van de
codetabel.
\end{exercise}

\begin{exercise}[$\star$]\label{exo:g11:gene-expression:4}
Wat zijn de rollen van het \omterm{def:g10:cells-common-unit:organelle}{ribosoom} en van de \omterm{def:g11:gene-expression:transcription}{transport-RNA}'s bij de
\omterm{prop:g11:gene-expression:translation}{translatie}?
\end{exercise}

\begin{exercise}[$\star$]\label{exo:g11:gene-expression:5}
Waarom moet de code minstens drie \omterm{def:g10:universal-dna:nucleotide}{nucleotiden} per \omterm{def:g10:chemistry-of-life:families}{aminozuur} gebruiken?
\end{exercise}

\begin{exercise}[$\star\star$]\label{exo:g11:gene-expression:6}
Een \omterm{def:g11:gene-expression:protein}{eiwit} telt 412 \omterm{def:g10:chemistry-of-life:families}{aminozuren}. Wat is de kleinste lengte van het
coderende deel van zijn mRNA, het stopcodon meegerekend, en van het
bijbehorende \omterm{def:g10:universal-dna:gene}{gen} in basenparen?
\end{exercise}

\begin{exercise}[$\star\star$]\label{exo:g11:gene-expression:7}
De coderende streng \texttt{ATG AAA GGC TGG TAA} wordt gemuteerd tot
\texttt{ATG AAA GGC TGA TAA}. Geef beide \omterm{def:g10:chemistry-of-life:families}{eiwitten} en deel de \omterm{def:g11:mutations:mutation}{mutatie} in.
\end{exercise}

\begin{exercise}[$\star\star$]\label{exo:g11:gene-expression:8}
Zelfde beginvolgorde, gemuteerd tot \texttt{ATG AAG GGC TGG TAA} en tot
\texttt{ATG AAA GGG CTG GTA A}. Deel elk in en geef de \omterm{def:g10:chemistry-of-life:families}{eiwitten}.
\end{exercise}

\begin{exercise}[$\star\star$]\label{exo:g11:gene-expression:9}
Leg met de tabel uit waarom een substitutie op de derde plaats van een
\omterm{def:g11:gene-expression:code}{codon} vaak stil is.
\end{exercise}

\begin{exercise}[$\star\star$]\label{exo:g11:gene-expression:10}
Poly-UC-RNA (\texttt{UCUCUCUC}\dots) geeft een \omterm{def:g11:gene-expression:protein}{eiwit} waarin serine en
leucine elkaar afwisselen. Laat zien dat dit strookt met een code van
drie letters en gebruik het om twee \omterm{def:g11:gene-expression:code}{codons} toe te wijzen.
\end{exercise}

\begin{exercise}[$\star\star$]\label{exo:g11:gene-expression:11}
Een menselijk \omterm{def:g10:universal-dna:gene}{gen} van \num{30000} basenparen brengt een \omterm{def:g11:gene-expression:protein}{eiwit} van 500
\omterm{def:g10:chemistry-of-life:families}{aminozuren} voort. Verklaar het verschil.
\end{exercise}

\begin{exercise}[$\star\star\star$]\label{exo:g11:gene-expression:12}
Een geneesmiddel blokkeert bacteriële \omterm{def:g10:cells-common-unit:organelle}{ribosomen} maar niet de menselijke.
Leg uit waarom het een antibioticum kan zijn, en wat zijn bestaan zegt
over de \omterm{def:g10:cells-common-unit:organelle}{ribosomen} van de twee groepen.
\end{exercise}

\begin{exercise}[$\star\star\star$]\label{exo:g11:gene-expression:13}
Leg in de woorden van dit hoofdstuk uit hoe het kwallengen uit
\cref{ch:g10:universal-dna} door een muis gelezen kon worden, en wat er
zou gebeuren als de muis een andere code gebruikte.
\end{exercise}

\begin{exercise}[$\star\star\star$]\label{exo:g11:gene-expression:14}
Een \omterm{def:g11:mutations:mutation}{mutatie} verandert het tRNA waarvan het anticodon UAG leest, zodat het
een \omterm{def:g10:chemistry-of-life:families}{aminozuur} draagt in plaats van te stoppen. Voorspel het gevolg voor
de \omterm{def:g10:chemistry-of-life:families}{eiwitten} van de \omterm{def:g10:cells-common-unit:cell}{cel} in het algemeen, en voor een gemuteerd \omterm{def:g10:universal-dna:gene}{gen} met een
voortijdige UAG.
\end{exercise}

\begin{exercise}[$\star\star\star$]\label{exo:g11:gene-expression:15}
Bespreek waarom een leesraamverschuiving vlak bij het begin van een \omterm{def:g10:universal-dna:gene}{gen}
vrijwel altijd erger is dan een vlak bij het eind ervan, en waarom een
\omterm{def:g11:mutations:mutation}{missense-mutatie} alles kan zijn van onschadelijk tot dodelijk.
\end{exercise}

\section{Probleem: Eén gen op vier manieren gelezen}

\begin{problem}[{Weekendprobleem --- één kort gen overgeschreven,
vertaald, drie keer gemuteerd en geklokt: van het DNA tot de aminozuren,
en de zestig seconden die een eiwit kost}]
\label{pb:g11:gene-expression:1}
De coderende streng van een kort \omterm{def:g10:universal-dna:gene}{gen} luidt:
\begin{center}
\texttt{ATG GCT TGG AAA CCC GAA TAC GGT CAT TTC TGA}
\end{center}

\medskip
\noindent\textbf{Deel I --- Het \omterm{def:g10:universal-dna:gene}{gen} lezen.}
\begin{enumerate}
  \item Schrijf de matrijsstreng op.
  \item Schrijf het mRNA op dat van het \omterm{def:g10:universal-dna:gene}{gen} wordt overgeschreven.
  \item Vertaal het \omterm{def:g11:gene-expression:code}{codon} voor \omterm{def:g11:gene-expression:code}{codon} en geef de aminozuurvolgorde van het
        \omterm{def:g11:gene-expression:protein}{eiwit}.
  \item Hoeveel \omterm{def:g10:chemistry-of-life:families}{aminozuren} telt het \omterm{def:g11:gene-expression:protein}{eiwit}, en hoeveel \omterm{def:g10:universal-dna:nucleotide}{nucleotiden} van het
        mRNA worden gebruikt om ze aan te wijzen, de stop meegerekend?
  \item Welk \omterm{def:g11:gene-expression:code}{codon} zette het lezen in gang, en wat legt het leesraam voor
        alle andere vast?
\end{enumerate}

\medskip
\noindent\textbf{Deel II --- Drie mutanten.}
\begin{enumerate}[resume]
  \item Mutant 1: \texttt{ATG GCT TGG AAA CCG GAA TAC GGT CAT TTC TGA}.
        Geef het \omterm{def:g11:gene-expression:protein}{eiwit} en deel de \omterm{def:g11:mutations:mutation}{mutatie} in.
  \item Mutant 2: \texttt{ATG GCT TGG AAA CCC GAA TAC GGT CAT TTC TGA}
        waarbij het negende \omterm{def:g11:gene-expression:code}{codon} \texttt{CAT} is vervangen door
        \texttt{CGT}. Geef het \omterm{def:g11:gene-expression:protein}{eiwit} en deel de \omterm{def:g11:mutations:mutation}{mutatie} in.
  \item Mutant 3: \texttt{ATG GCT TGA AAA CCC GAA TAC GGT CAT TTC TGA}.
        Geef het \omterm{def:g11:gene-expression:protein}{eiwit} en deel de \omterm{def:g11:mutations:mutation}{mutatie} in. Welke enkele substitutie
        bracht haar voort?
  \item Mutant 4: er wordt na de zesde \omterm{def:g10:universal-dna:nucleotide}{nucleotide} een T ingevoegd:
        \texttt{ATG GCT TTG GAA ACC CGA ATA CGG TCA TTT CTG A}. Geef het
        \omterm{def:g11:gene-expression:protein}{eiwit} en deel de \omterm{def:g11:mutations:mutation}{mutatie} in.
  \item Zet de vier mutanten op volgorde van minst tot meest ontwrichtend
        voor het \omterm{def:g11:gene-expression:protein}{eiwit} en verantwoord dat.
\end{enumerate}

\medskip
\noindent\textbf{Deel III --- Waarom de code is zoals zij is.}
\begin{enumerate}[resume]
  \item Hoeveel \omterm{def:g11:gene-expression:code}{codons} wijzen leucine aan? En serine? En tryptofaan?
        Welk \omterm{def:g10:chemistry-of-life:families}{aminozuur} wordt het waarschijnlijkst door een willekeurige
        substitutie in zijn \omterm{def:g11:gene-expression:code}{codon} veranderd, en welk het minst?
  \item Bereken welk deel van alle substituties op de derde plaats van
        een leucinecodon \texttt{CUx} stil is.
  \item Leg uit waarom een code van twee letters niet zou kunnen werken,
        en waarom een code van vier niet nodig is.
  \item Poly-U geeft poly-fenylalanine en poly-A poly-lysine. Welke
        vermeldingen in de tabel stellen die twee proeven vast?
  \item De code is dezelfde bij de bacterie en bij de mens. Geef het
        gevolg voor de \omterm{prop:g10:universal-dna:universal}{transgenese} en het gevolg voor de verwantschap.
\end{enumerate}

\medskip
\noindent\textbf{Deel IV --- De fabriek klokken.}
Een \omterm{def:g10:cells-common-unit:organelle}{ribosoom} voegt 10 \omterm{def:g10:chemistry-of-life:families}{aminozuren} per seconde toe en \omterm{def:g10:cells-common-unit:organelle}{ribosomen} volgen
elkaar langs een mRNA met 80 \omterm{def:g10:universal-dna:nucleotide}{nucleotiden} ertussen; elk mRNA leeft 2 uur.
\begin{enumerate}[resume]
  \item Hoelang doet één \omterm{def:g10:cells-common-unit:organelle}{ribosoom} over het vertalen van het \omterm{def:g11:gene-expression:protein}{eiwit} van dit
        \omterm{def:g10:universal-dna:gene}{gen}? En over een \omterm{def:g11:gene-expression:protein}{eiwit} van 600 \omterm{def:g10:chemistry-of-life:families}{aminozuren}?
  \item Hoeveel \omterm{def:g10:cells-common-unit:organelle}{ribosomen} kunnen een mRNA van 1800 \omterm{def:g10:universal-dna:nucleotide}{nucleotiden} tegelijk
        lezen?
  \item Hoeveel kopieën van het \omterm{def:g11:gene-expression:protein}{eiwit} levert één mRNA in zijn leven op,
        als er elke 8 seconden een nieuw \omterm{def:g10:cells-common-unit:organelle}{ribosoom} begint?
  \item Een bacterie heeft binnen 10 minuten \num{2000} kopieën van een
        \omterm{def:g10:cell-metabolism:metabolism}{enzym} nodig. Hoeveel mRNA's van het \omterm{def:g10:universal-dna:gene}{gen}, tegelijk overgeschreven,
        zijn daarvoor nodig?
  \item Geef het resultaat: het \omterm{def:g11:gene-expression:protein}{eiwit} waarvoor het \omterm{def:g10:universal-dna:gene}{gen} codeert, de
        \omterm{def:g11:mutations:mutation}{mutatie} van de vier die het teniet deed, en de tijd die één
        \omterm{def:g10:cells-common-unit:organelle}{ribosoom} nodig heeft om het te bouwen.
\end{enumerate}
\end{problem}
